\documentclass[11pt]{article}
\input{../../shared/project_style.tex}
\rhead{Basics 44 Textbook Draft}

\begin{document}
\DocTitle{Basics 44: Guiding-Center Theory}{VI - Kinetic theory and energetic particles}{Separation of gyromotion from slow orbit dynamics, adiabatic invariants, mirror force, and reduced equations of motion}

\begin{overviewbox}
\textbf{Textbook draft, version 1.0.} A fusion alpha particle can execute billions of gyro-orbits before its macroscopic confinement changes. Guiding-center theory removes that rapidly rotating phase while preserving the slow motion needed for orbit following, resonances, and wall-loss calculations. This chapter develops the ordering, magnetic moment, mirror force, principal drifts, validity conditions, and numerical implications of the reduction.
\end{overviewbox}

\ProjectTOC

\section{Learning objectives}
After completing this note, the reader should be able to:
\begin{itemize}[leftmargin=1.6em]
\item Identify the small parameter that justifies gyro-averaging.
\item Separate particle position into guiding-center and gyroradius contributions.
\item Derive and interpret the first adiabatic invariant $\mu$.
\item Derive the mirror force and trapped-particle condition qualitatively.
\item State first-order guiding-center equations and principal drifts.
\item Recognize when guiding-center theory fails.
\end{itemize}

\section{Prerequisites and reading strategy}
Recommended prerequisites are Basics 2, 13, 41, and 42. Review uniform-field gyromotion before studying nonuniform-field corrections.

\section{Separation of scales}
In a magnetic field, a charged particle rotates at gyrofrequency
\begin{equation}
\Omega_c=\frac{|q|B}{m}
\end{equation}
with Larmor radius
\begin{equation}
\rho_L=\frac{v_\perp}{\Omega_c}.
\end{equation}
If the equilibrium varies over length $L_B$ and time $\tau_B$, guiding-center theory requires
\begin{equation}
\epsilon=\frac{\rho_L}{L_B}\ll1,
\qquad
\frac{1}{\Omega_c\tau_B}\ll1.
\end{equation}
The particle position is decomposed as
\begin{equation}
\mathbf r=\mathbf R+\boldsymbol\rho,
\end{equation}
where $\mathbf R$ is the slowly moving guiding center and $\boldsymbol\rho$ is the rapidly rotating gyroradius vector.

\section{Uniform-field orbit}
For $\mathbf E=0$ and uniform $\mathbf B=B\mathbf b$,
\begin{equation}
m\dot{\mathbf v}=q\mathbf v\times\mathbf B.
\end{equation}
The parallel velocity is constant, while the perpendicular velocity rotates. The kinetic energy and pitch components satisfy
\begin{equation}
\frac{mv^2}{2}=\frac{mv_\parallel^2}{2}+\frac{mv_\perp^2}{2}.
\end{equation}
This exact solution is the zeroth-order motion about which the asymptotic expansion is built.

\section{Magnetic moment as an adiabatic invariant}
Define
\begin{equation}
\boxed{\mu=\frac{mv_\perp^2}{2B}}.
\end{equation}
The perpendicular kinetic energy is $E_\perp=\mu B$. When $B$ changes little during one gyro-orbit, gyro-averaging shows
\begin{equation}
\frac{d\mu}{dt}=O(\epsilon^2)
\end{equation}
for slowly varying static fields. Thus $\mu$ is approximately conserved.

A useful physical derivation treats the gyro-orbit as a current loop. Its magnetic moment is
\begin{equation}
I\mathcal A=\frac{q\Omega_c}{2\pi}\pi\rho_L^2
=\frac{mv_\perp^2}{2B}\operatorname{sgn}(q).
\end{equation}
The orbital magnetic dipole interacts with field gradients, producing a force proportional to $-\mu\nabla B$.

\section{Mirror force}
For static fields and negligible electric work, the total kinetic energy
\begin{equation}
\mathcal E=\frac{mv_\parallel^2}{2}+\mu B
\end{equation}
is constant. Differentiate along the orbit:
\begin{equation}
m v_\parallel\dot v_\parallel+\mu v_\parallel\mathbf b\cdot\nabla B=0.
\end{equation}
For $v_\parallel\neq0$,
\begin{equation}
\boxed{m\dot v_\parallel=-\mu\nabla_\parallel B}.
\end{equation}
A particle moving into larger $B$ converts parallel energy into perpendicular energy until $v_\parallel=0$, where it reflects.

If the particle begins where the field is $B_0$ with pitch angle $\alpha_0$, reflection occurs when
\begin{equation}
\sin^2\alpha_0\ge\frac{B_0}{B_{\max}}.
\end{equation}
This divides trapped and passing populations.

\section{Guiding-center drifts}
The lowest-order guiding-center velocity is parallel streaming plus the electric drift:
\begin{equation}
\dot{\mathbf R}=v_\parallel\mathbf b+
\mathbf v_E+\mathbf v_{\nabla B}+\mathbf v_c+\cdots.
\end{equation}
The principal drifts are
\begin{align}
\mathbf v_E&=\frac{\mathbf E\times\mathbf B}{B^2},\\
\mathbf v_{\nabla B}&=\frac{\mu}{qB}\,\mathbf b\times\nabla B,\\
\mathbf v_c&=\frac{mv_\parallel^2}{qB}\,\mathbf b\times\boldsymbol\kappa,
\qquad
\boldsymbol\kappa=\mathbf b\cdot\nabla\mathbf b.
\end{align}
The $\mathbf E\times\mathbf B$ drift is independent of particle mass and charge. Grad-$B$ and curvature drifts reverse with charge and grow with particle energy, making them central to fast-ion confinement.

A consistent first-order parallel equation is
\begin{equation}
\boxed{\dot v_\parallel=\frac{q}{m}E_\parallel-
\frac{\mu}{m}\mathbf b\cdot\nabla B+\cdots.}
\end{equation}

\section{Guiding-center coordinates and measure}
A common coordinate set is
\begin{equation}
(\mathbf R,v_\parallel,\mu,\vartheta_g),
\end{equation}
where $\vartheta_g$ is gyroangle. Gyro-averaged models remove explicit dependence on $\vartheta_g$. The phase-space measure is not simply $d^3R\,dv_\parallel\,d\mu$; a magnetic Jacobian factor appears. Structure-preserving formulations use a modified field $B_\parallel^*$ and measure proportional to
\begin{equation}
B_\parallel^*\,d^3R\,dv_\parallel\,d\mu\,d\vartheta_g.
\end{equation}
Ignoring this factor can violate particle conservation in noncanonical coordinates.

\section{Orbit frequencies and resonance}
A guiding-center orbit can contain fast parallel transit or bounce motion and slower toroidal precession. These generate frequencies
\begin{equation}
\omega_t,\qquad\omega_b,\qquad\omega_d,
\end{equation}
which enter Alfv\'en-wave resonance conditions. Module 3 must compute these consistently with the same equilibrium and coordinate conventions used by Module 4.

\section{Energy and invariant checks}
In time-independent electrostatic fields, an ideal guiding-center integrator should preserve
\begin{equation}
\mathcal H=\frac{mv_\parallel^2}{2}+\mu B+q\phi
\end{equation}
to numerical accuracy. In an axisymmetric field, canonical toroidal momentum provides another benchmark. Monitoring only spatial trajectories can hide large invariant errors.

\section{Limits of validity}
Guiding-center theory can fail when:
\begin{itemize}[leftmargin=1.6em]
\item $\rho_L/L_B$ is not small;
\item fields vary near the gyrofrequency;
\item the orbit passes close to magnetic nulls where $B\to0$;
\item sharp material boundaries require full-orbit impact geometry;
\item cyclotron resonances or finite-Larmor-radius effects are essential.
\end{itemize}
Near marginal validity, a full-orbit comparison is required.

\section{Numerical considerations}
Runge--Kutta integration is simple but can accumulate secular invariant drift. Symplectic, variational, or volume-preserving methods better respect long-time orbit structure. Adaptive stepping must resolve bounce points and rapid field variation without introducing biased loss boundaries.

\section{Connection to the Kinetic MHD-PINN project}
Module 3 uses guiding-center trajectories to estimate confinement, orbit frequencies, prompt losses, and wall strikes. Module 5 uses orbit frequencies and mode overlap for resonant drive. Structure-preserving neural surrogates may accelerate orbit maps, but they must be tested against energy, magnetic-moment, and phase-space-volume constraints.

\section{Computational laboratory}
\begin{enumerate}[leftmargin=1.6em]
\item Integrate full Lorentz motion in a uniform field and recover $\Omega_c$ and $\rho_L$.
\item Add a weak field gradient and compare the measured grad-$B$ drift with theory.
\item Integrate the reduced mirror equations in $B(z)=B_0(1+z^2/L^2)$.
\item Verify conservation of $\mu$ and energy as the scale ratio $\rho_L/L$ is reduced.
\item Compare ordinary Runge--Kutta and a structure-preserving method over many bounce periods.
\end{enumerate}

\section{Summary}
Guiding-center theory exploits the small ratio of gyroradius to equilibrium scale. It removes fast gyroangle dynamics while retaining parallel streaming, mirror motion, and cross-field drifts. The magnetic moment is the key first adiabatic invariant. The reduction enables practical energetic-particle orbit following, but its Jacobian, invariants, and validity limits must be respected.

\section{Exercises}
\begin{enumerate}[leftmargin=1.6em]
\item Estimate the gyroradius of a 3.5 MeV alpha particle in a 5 T field.
\item Derive the mirror reflection condition from energy and $\mu$ conservation.
\item Show dimensionally that each drift velocity has units of speed.
\item Explain why $\mathbf E\times\mathbf B$ drift is independent of charge sign.
\item Propose invariant-based acceptance tests for a Module 3 orbit solver.
\end{enumerate}

\SymbolsAcronymsSection
\begin{symtable}
$\Omega_c$ & Gyrofrequency & Angular frequency of cyclotron motion.\\
$\rho_L$ & Larmor radius & Radius of the fast gyromotion.\\
$\mu$ & Magnetic moment & First adiabatic invariant $mv_\perp^2/(2B)$.\\
$\mathbf R$ & Guiding-center position & Gyro-averaged particle position.\\
$\boldsymbol\kappa$ & Field-line curvature & $\mathbf b\cdot\nabla\mathbf b$.\\
\end{symtable}

\section{References}
\begin{enumerate}[leftmargin=1.6em]
\item R. G. Littlejohn, guiding-center Hamiltonian theory papers.
\item R. B. White, \emph{The Theory of Toroidally Confined Plasmas}.
\item T. G. Northrop, \emph{The Adiabatic Motion of Charged Particles}.
\item A. J. Brizard and T. S. Hahm, ``Foundations of nonlinear gyrokinetic theory,'' \emph{Reviews of Modern Physics}, 2007.
\end{enumerate}

\end{document}
